\section{Lyndon's original homomorphic-image preservation theorem}
\label{sec:p03-n-a-011-02}

Lyndon's theorem concerns \emph{images}: the homomorphisms in this
presentation are onto. It admits operation symbols, including constants,
as well as relations and equality. The stronger result describes the
positive theory of an elementary class by elementary extensions of its
images; sentence preservation is a consequence, not its starting point
\cite[pp.~143--154]{Lyndon1959}.

Here ``elementary class'' permits a set of axioms, not only one sentence.
Lyndon places the paper inside Tarski's theory of elementary classes: he
prefers ``elementary'' to ``arithmetical'', so that his elementary classes
are Tarski's arithmetical classes in the wider sense, written
\(AC_{\Delta}\), and he credits the notion of elementary extension to
Tarski and Vaught \cite[p.~144 and n.~1]{Lyndon1959}.

Lyndon contrasts the two directions of the sentence theorem. That every
sentence equivalent to a positive one is preserved he describes as not
hard to show; the converse, although it seems almost as evident
intuitively, he describes as a matter of considerable difficulty
\cite[p.~143]{Lyndon1959}. Preservation of positive sentences had been
observed by Lyndon himself and by Marczewski, who asked whether the
converse holds. \L o\'s had announced a proof using methods quite unlike
Lyndon's, which had not appeared in print; Mal'cev had stated the result,
and Lyndon says only that Mal'cev's paper \emph{appears} to indicate a
method of proof \cite[p.~143]{Lyndon1959}. Lyndon thanks Henkin and Tarski
for many discussions, held in particular during his visit to the University
of California at Berkeley, and records that Robinson suggested the
relativization of Theorem~5\('\) at the American Mathematical Society
Summer Institute in Ithaca in 1957 \cite[p.~144 and nn.~2--3]{Lyndon1959}.
These are the source's historical attributions, not alternative proofs
supplied by the present article \cite[pp.~143--144,151]{Lyndon1959}.

\subsection{Language, polarity, and the imported interpolation result}

\begin{hbconvention}{def:p03-lyndon-language}{The auxiliary language and genuine identity}
A language has an infinite set \(V\) of variables, operation symbols
\(W\), relation symbols \(R\), and a finite arity \(\rho(s)\) for each
\(s\in W\cup R\). Operations are total, and domains are nonempty.
Terms and atomic formulas have the forms \(w(t_1,\ldots,t_{\rho(w)})\)
and \(r(t_1,\ldots,t_{\rho(r)})\). The logical primitives are
\(0,1,\neg,\land,\lor,\forall,\exists\), with \(0\) false and
\(1\) true. Empty conjunctions and disjunctions are \(1\) and \(0\).

Initially every relation symbol is freely interpreted. In particular,
the distinguished binary symbol \(e\), when present, is not yet logical
identity. A \emph{relational system} is a structure in which \(e\)
is the actual identity relation. Congruence axioms and quotients connect
the auxiliary structures to relational systems; replacing auxiliary
\(e\) by logical equality prematurely changes the argument.
\end{hbconvention}

\begin{hbdefinition}{def:p03-lyndon-polarity}{Polarity and increasing theories}
An occurrence of a relation symbol is positive or negative according as
the number of enclosing negations is even or odd. A formula is
\(Q\)-positive, for \(Q\subseteq R\), if no member of \(Q\) has
a negative occurrence. Negations of atoms outside \(Q\) are permitted.
A positive formula is \(R\)-positive, equivalently, after normalization,
one built from atoms and \(0,1\) by \(\land,\lor,\forall,\exists\).

For distinct fresh symbols \(q'\) of the same arities as \(q\in Q\),
write \(F'\) for simultaneous replacement of these symbols, and put
\[
 I_Q=\{\forall\bar x\,(q(\bar x)\to q'(\bar x)):q\in Q\}.
\]
A set \(\Gamma\) of formulas is \emph{increasing in \(Q\)} if
\(\Gamma\cup I_Q\entails\Gamma'\), for every assignment when formulas
have free variables. Thus functions, domain, and relations outside \(Q\)
remain fixed while the selected relations increase.
\end{hbdefinition}

\begin{hbtheorem}{thm:p03-lyndon-interpolation-import}{cited=Lyndon1959Interpolation}{The precise polarity interpolation dependency}
In the auxiliary predicate calculus, if sentences \(S,T\) of \(L\) satisfy
\(\{S\}\entails T\), there is a sentence \(M\) of the same language such that
\(\{S\}\entails M\) and \(\{M\}\entails T\). A relation symbol
has a positive, respectively negative, occurrence in \(M\) only if
it has an occurrence of that same polarity in \emph{both} \(S\) and
\(T\). This is a restriction on relation occurrences, not a claim in
this formulation that every operation symbol occurs on both sides.
\end{hbtheorem}

This is the theorem explicitly imported on preservation p.~144 from the
distinct article \emph{An Interpolation Theorem in the Predicate Calculus},
pp.~129--142. That article proves it through Skolemization, duplication,
specialization, and a propositional inconsistency theorem, rather than
through a sequent calculus \cite[pp.~133--141]{Lyndon1959Interpolation}.
Its final identity version includes \(e\) among the symbols subject to
the polarity condition, with equality closure added to the Boolean-ideal
and literal arguments. The preservation proof below uses the freely
interpreted auxiliary version and only subsequently takes quotients.

\begin{hblemma}{thm:p03-lyndon-increasing}{proved=annex:p3.n-a-011.2.increasing}{Propositions 1 and 2 and Corollary 2.1}
Every \(Q\)-positive formula is increasing in \(Q\). More generally,
let \(\Sigma,\Gamma,\Delta\) be sets of sentences in \(L\). If
\[
 \Sigma\cup\Sigma'\cup\Gamma\cup I_Q\entails\Delta',
\]
there is a set \(\Pi\) of \(Q\)-positive sentences, with no symbols
from \(Q'\), such that
\[
 \Sigma\cup\Gamma\entails\Pi,
 \qquad \Sigma\cup\Pi\entails\Delta.
\]
Consequently a set of sentences is increasing in \(Q\) exactly when it
is equivalent to a set of \(Q\)-positive sentences.
\end{hblemma}

\subsection{Surjective \(Q\)-maps and the elementary diagram}

\begin{hbdefinition}{def:p03-lyndon-qmap}{The source's \(Q\)-maps and closure notation}
A \(Q\)-map \(h:A\to B\) is a surjection satisfying, for every tuple,
\[
 h(w^A(\bar a))=w^B(h\bar a),\qquad
 r^A(\bar a)\Rightarrow r^B(h\bar a),
\]
where the latter implication is an equivalence if \(r\notin Q\).
A bijective \(Q\)-map is a \(Q\)-isomorphism. In particular,
\(\varnothing\)-maps preserve and reflect every atomic relation;
their bijective instances are ordinary isomorphisms.

For a set \(\Gamma\) of sentences, \(\Gamma^*\) is its model class.
For a class \(K\), \(K^*\) is its common sentence theory. Hence
\(\Gamma^{**}\) is logical closure and \(K^{**}\) is elementary
closure. An elementary class has the form \(\Gamma^*\).
An elementary extension \(B\) of \(A\) contains \(A\) as a substructure
and agrees on every formula with every assignment from \(A\).
\end{hbdefinition}

\begin{hblemma}{thm:p03-lyndon-map-calculus}{proved=annex:p3.n-a-011.2.maps}{Terms, formulas, kernels, and quotients}
\(Q\)-maps preserve truth of \(Q\)-positive formulas with assignments
transported along the map. A \(\varnothing\)-map preserves and reflects
every formula. Its kernel is a congruence for all operations and relations;
it factors through the congruence quotient as an isomorphism. Conversely
the quotient map by any congruence is a \(\varnothing\)-map. These claims
apply to the auxiliary language, before \(e\) is required to be identity.
\end{hblemma}

\begin{hbdefinition}{def:p03-lyndon-diagram}{The expanded diagram used for Proposition~3}
For \(A\), let \(L_A\) add constants \(w_a\), one for each \(a\in A\),
and a fresh binary relation \(e_A\). In \(A_A\) the constants denote
their indicated elements and \(e_A\) is actual identity. Its complete
sentence theory \(A_A^*\), rather than only its atomic diagram, is used
below. The new \(e_A\) remains a freely interpreted relation in other
\(L_A\)-structures.
\end{hbdefinition}

Lyndon does not claim this apparatus as new. He notes that his \(Q\)-maps
and the notion of elementary extension are special cases of more general
ideas then being developed by Keisler, that the constants \(w_a\) come from
Robinson's diagrams, and that Proposition~3 is contained in Theorem~1.11 of
Tarski and Vaught \cite[p.~146, n.~4]{Lyndon1959}.

\begin{hblemma}{thm:p03-lyndon-diagram}{proved=annex:p3.n-a-011.2.diagram}{Proposition 3: a diagram followed by a quotient}
If an \(L\)-structure \(B\) has an \(L_A\)-expansion satisfying
\(A_A^*\), then a \(\varnothing\)-image of \(B\), after relabeling
the named copy of \(A\), is an elementary extension of \(A\).
\end{hblemma}

\begin{hbtheorem}{thm:p03-lyndon-qimage}{proved=annex:p3.n-a-011.2.qimage}{Theorem 4: the \(Q\)-positive theory of an elementary class}
For any auxiliary \(L\)-structure \(A\), elementary class \(K\),
and \(Q\subseteq R\), the following are equivalent:
\begin{enumerate}
\item \(A\) satisfies every \(Q\)-positive sentence in \(K^*\).
\item Some elementary extension of \(A\) is a \(Q\)-image of a member
      of \(K\).
\end{enumerate}
\end{hbtheorem}

\begin{hbcorollary}{thm:p03-lyndon-qclosure}{proved=annex:p3.n-a-011.2.qclosure}{Corollaries 4.1 and 4.2}
An elementary class is closed under \(Q\)-images exactly when it is
axiomatizable by \(Q\)-positive sentences. Equivalently, a set of
sentences is preserved under \(Q\)-maps exactly when it is equivalent
to a set of \(Q\)-positive sentences.
\end{hbcorollary}

\subsection{Quotienting to relational systems}

\begin{hbdefinition}{def:p03-lyndon-image-operations}{Homomorphism, projection, and enlargement}
In relational systems a homomorphism is an onto map preserving all
operations and relations: an \(R\)-map. A \emph{projection} is a
quotient by a congruence for operations and all relations except logical
identity; a \emph{relation enlargement} is a bijection commuting with
operations and increasing the relations other than \(e\).

The auxiliary \(\varnothing\)-map notion must not be substituted for
a possibly noninjective projection between genuine-identity systems:
reflection of identity would make it injective. At the relational-system
stage the projection is a \(\{e\}\)-map, and an enlargement is an
\(R\setminus\{e\}\)-map. Every onto homomorphism factors as an
enlargement followed by a projection: enlarge each nonequality relation
of its source to the inverse image of the target relation.
\end{hbdefinition}

\begin{hbtheorem}{thm:p03-lyndon-main}{proved=annex:p3.n-a-011.2.main}{Theorem 5: three relational-system cases}
Let \(A\) be a relational system and \(K\) an elementary class of
relational systems. The following pairs are equivalent, row by row:
\begin{center}
\begin{tabular}{p{0.46\textwidth}p{0.46\textwidth}}
\(A\) satisfies the sentences of \(K^*\) positive in &
An elementary extension of \(A\) is a member of \\
\hline
all relation symbols & the homomorphic images of \(K\);\\
\(e\) & the projections of \(K\);\\
all relation symbols except \(e\) & the enlargements of \(K\).
\end{tabular}
\end{center}
\end{hbtheorem}

The source also notes that the same argument can accommodate a selected
\(Q\) containing or excluding \(e\), and maps increasing one set of
relations while decreasing another. It does not separately develop those
mixed-direction formulations, choosing the homomorphism case for the
remaining corollaries \cite[p.~150]{Lyndon1959}.

\begin{hbcorollary}{thm:p03-lyndon-positive-corollaries}{proved=annex:p3.n-a-011.2.positive}{Corollaries 5.1--5.3, including sentence preservation}
For an elementary class \(K\) of relational systems, a sentence \(S\)
holds in every homomorphic image of \(K\) exactly when
\(\{P\}\entails S\) for some positive sentence \(P\) true throughout
\(K\). The class \(K\) is closed under homomorphic images exactly
when it is axiomatizable, among relational systems, by positive sentences.
In particular, a sentence is preserved under onto homomorphisms between
all relational systems exactly when it is equivalent to a positive sentence.
\end{hbcorollary}

\begin{hbtheorem}{thm:p03-lyndon-two-extensions}{proved=annex:p3.n-a-011.2.two-extensions}{The strengthened two-system observation}
\(A\) satisfies every positive sentence true in a relational system
\(B\) exactly when there are elementary extensions \(A'\) of \(A\)
and \(B'\) of \(B\) and an onto homomorphism \(B'\to A'\).
\end{hbtheorem}

\begin{hbtheorem}{thm:p03-lyndon-relative}{proved=annex:p3.n-a-011.2.relative}{Theorem 5' and its relative corollaries}
Let \(H\) be an elementary class of relational systems, \(A\in H\),
and \(K\) an elementary subclass of \(H\). Then \(A\) satisfies the
positive sentences true in \(K\) exactly when an elementary extension
of \(A\) is a homomorphic image of a member of \(K\).

All three positive corollaries hold relative to \(H\): replace consequence
and equivalence by consequence and equivalence among \(H\)-systems,
and consider only images that also lie in \(H\). In particular a sentence
\(S\) is preserved by onto homomorphisms between \(H\)-systems exactly
when it is \(H\)-equivalent to a positive sentence.
\end{hbtheorem}

For torsionfree groups, \(H\) is axiomatized by the group axioms and
\(\forall x\,(x^n=1\to x=1)\), \(n\geq1\). The relative result
does not say every quotient of a torsionfree group is torsionfree. It
compares only maps whose two endpoints satisfy those axioms, and says the
equivalence to a positive sentence follows from those same axioms
\cite[p.~151]{Lyndon1959}.

\subsection{Examples and limits of the theorem}

\begin{hbtheorem}{thm:p03-lyndon-algebra-examples}{proved=annex:p3.n-a-011.2.algebra}{Commutativity and cancellation}
Commutativity \(\forall x\forall y\,xy=yx\) is positive and preserved
under onto homomorphisms. Left cancellation
\(\forall x\forall y\forall z\,(xy=xz\to y=z)\) is not preserved
under such homomorphisms, and is not equivalent over all multiplicative
systems to a positive sentence.
\end{hbtheorem}

\begin{hbtheorem}{thm:p03-lyndon-nonelementary}{proved=annex:p3.n-a-011.2.nonelementary}{Why the source class must be elementary}
In pure identity language every set satisfies the positive sentences true
in all nonempty finite sets. Nevertheless an infinite set has no elementary
extension that is an onto image of a finite set. Thus elementary closure
cannot be omitted from the hypothesis on \(K\).
\end{hbtheorem}

\begin{hbdefinition}{def:p03-lyndon-chain-example}{The complementary chain sentences}
In language \(\{r,e\}\), with \(r\) binary and \(e\) genuine identity,
put
\[
 S=\exists x\exists y\forall z\exists t\,
   \bigl(r(x,y)\land
       (r(x,z)\to(r(x,t)\land r(z,t)))\bigr)
\]
and, for \(n\geq1\),
\[
 S_n=\exists x_0\cdots x_n\,
 \left(\bigwedge_{i=1}^{n}r(x_0,x_i)
       \land\bigwedge_{i=1}^{n-1}r(x_i,x_{i+1})\right).
\]
The initial edge in \(S\) is a conjunct outside the implication. No
pairwise-distinctness conditions occur in \(S_n\).
\end{hbdefinition}

\begin{hbdefinition}{def:p03-lyndon-constructive-operations}{Skolemization, duplication, and specialization in the imported theorem}
Write a prenex sentence, allowing empty universal blocks, as
\[
 S=\forall\bar x_1\exists y_1\cdots
       \forall\bar x_n\exists y_n\,M,
\]
where the quantified variables are distinct and \(M\) has no quantifiers.
Its Skolem substitution \(\sigma\) fixes the universal variables and
replaces \(y_r\) by a fresh operation
\(s_r(\bar x_1,\ldots,\bar x_r)\). Its Skolem form universally closes
\(\sigma M\). A model satisfies \(S\) iff some expansion satisfies
its Skolem form; the expansion chooses witnesses as functions of the
preceding universal variables.

A \emph{duplication} takes finitely many variable-disjoint renamed copies
\(\pi_iM\), conjoins them, and quantifies their variables in an order
such that \(\forall x_{pq}^{i}\) precedes \(\exists y_r^{i}\) whenever
\(p\leq r\). A \emph{specialization} fixes every \(y_r\) and replaces
each \(x_{pq}\) by a term in fresh universal variables \(u_h\) and
the existential variables \(y_s\) with \(s<p\). Its prefix quantifies
all these variables in an order subject to the following dependencies:
\(\forall u_h\) precedes \(\exists y_r\) if \(u_h\) occurs in any
replacement for \(x_{pq}\) with \(p\leq r\); and \(\exists y_s\)
precedes \(\exists y_r\) if \(y_s\) occurs in such a replacement.
The source's Propositions~6 and 7 prove that both operations give a
semantic consequence of the original sentence.

The two Skolemized matrices are \emph{propositionally inconsistent}
when a substitution for individual variables into terms makes their
conjunction a Boolean contradiction and is one-to-one on the atomic
formulas occurring in each matrix separately. This injectivity does not
require the two matrices' atom sets to have disjoint images: opposed
literals from the two sides must be allowed to meet. These are the
operations and conditions of the imported source theorem
\cite[pp.~134--137]{Lyndon1959Interpolation}.
\end{hbdefinition}

\begin{hbtheorem}{thm:p03-lyndon-constructive-import}{cited=Lyndon1959Interpolation}{The constructive dependency used in the chain example}
For inconsistent prenex sentences \(S^1,S^2\), the interpolation article's
Main Theorem gives prenex \(T^1,T^2,U^1,U^2\): each \(T^\delta\)
arises from \(S^\delta\) by duplication; each \(U^\delta\) arises
from \(T^\delta\) by specialization; and their Skolem matrices become
propositionally inconsistent under a substitution injective on the atoms
of each matrix. Its interpolation construction deletes literals from the
first matrix to retain only polarities opposed by the second matrix.
\end{hbtheorem}

\begin{hbtheorem}{thm:p03-lyndon-image-class-boundary}{proved=annex:p3.n-a-011.2.chain}{Why elementary extensions cannot be deleted}
For \(K=\{A:A\entails S\}\), the class \(H(K)\) of onto homomorphic
images need not be elementary. Its elementary closure is
\[
 H(K)^{**}=\{A:A\entails S_n\text{ for every }n\geq1\}.
\]
More precisely, if \(P\) is positive and \(\{S\}\entails P\),
then \(\{S_n\}\entails P\) for some \(n\geq1\).
The structure on \(\mathbb N\) with \(r(i,j)\) iff \(j<i\)
belongs to this elementary closure but not to \(H(K)\).
\end{hbtheorem}

\begin{hbtheorem}{thm:p03-lyndon-undecidable}{proved=annex:p3.n-a-011.2.undecidable}{The final undecidability reduction}
Let \(T\) be a consistent finitely axiomatizable undecidable theory in
an effective first-order language. Suppose every \(T\)-system has an
onto homomorphic image that is a one-element \(T\)-system. Then deciding
whether a sentence is \(T\)-equivalent to a positive sentence is impossible.
Lyndon's examples use undecidable pure predicate validity with a relation
of arity greater than one, and the elementary theory of groups; their
undecidability is an external premise of the reduction, not proved in this
paper. For these premises Lyndon refers to Tarski, Mostowski and Robinson's
\emph{Undecidable Theories} for the main concepts, to Church for the
undecidability of the predicate calculus, and, for the elementary theory of
groups, to the same monograph together with the reference to Tarski given
there and to Novikov's work on the word problem
\cite[p.~153, nn.~5--7]{Lyndon1959}.
\end{hbtheorem}

\begin{hbtheorem}{thm:p03-lyndon-lookup-boundaries}{proved=annex:p3.n-a-011.2.boundaries}{The operation and syntax boundary}
Existential-positive formulas are preserved by arbitrary homomorphisms;
positive formulas with universal quantifiers are preserved by onto
homomorphisms. The latter statement cannot replace ``onto'' by
``arbitrary.'' These are all-structure results, not finite-only preservation
results. The original sentence theorem extends to formulas with free
variables by naming the tuple with constants, preserving their values under
the map.
\end{hbtheorem}

Thus \(\forall x\,P(x)\) is positive but not preserved by the inclusion
of a one-point all-\(P\) structure into a two-point structure with just
one \(P\)-point. The false sentence \(0\) must remain in the positive
fragment: without it, in pure identity language all positive sentences
are true in a one-element system, so an inconsistent sentence would have
no positive equivalent. None of the compactness constructions above keeps
the resulting extension finite.
